\documentclass[../report_main.tex]{subfiles}
\begin{document}

\begin{lstlisting}
// Paste your code here. E.g.

package com.phungminhphuc;

import javafx.application.Application;
import javafx.scene.Scene;
import javafx.scene.control.Button;
import javafx.scene.layout.StackPane;
import javafx.stage.Stage;
public class JavaFXHello extends Application {
    private Button btnHello; // Declare a "Button" control
    
    @Override
    public void start(Stage primaryStage) {
        // Construct the "Button" and attach an "EventHandler"
        btnHello = new Button();
        btnHello.setText("Say Hello");
        btnHello.setOnAction(evt -> System.out.println("Hello World!"));
        // Construct a scene graph of nodes
        StackPane root = new StackPane(); 
        root.getChildren().add(btnHello); 
        Scene scene = new Scene(root, 500, 200); 
        primaryStage.setScene(scene); // The stage sets scene
        primaryStage.setTitle("Hello"); // Set window's title
        primaryStage.show(); // Set visible (show it)
    }
    public static void main(String[] args) {
        launch(args);
    }
}
\end{lstlisting}


\begin{sidewaysfigure}
  \centering
  \includegraphics[width=\textwidth]{Figures/ClassDiagramFinal-Core Model.jpg}
  \caption{Core Model Class Diagram}
\end{sidewaysfigure}

\begin{figure}[H]
  \centering
  \includegraphics[width=\textwidth]{Figures/ClassDiagramFinal-UI Rendering Handlers.jpg}
  \caption{UI, Rendering, and Handler Class Diagram}
\end{figure}

\begin{sidewaysfigure}
  \centering
  \includegraphics[width=\textwidth]{Figures/ClassDiagramFinal-Vehicle Behavior.jpg}
  \caption{Vehicle and Driver Behavior Class Diagram}
\end{sidewaysfigure}


\end{document}